\documentclass[11pt, a4paper]{article}
\usepackage[utf8]{inputenc}
\usepackage{amsmath,setspace,geometry}
\usepackage{amsthm}
\usepackage{amsfonts}
\usepackage{mathtools}
\mathtoolsset{showonlyrefs}
\usepackage[shortlabels]{enumitem}
\usepackage{rotating}
\usepackage{pdflscape}
\usepackage{graphicx}
\usepackage{bbm}
\usepackage[dvipsnames]{xcolor}
\usepackage[colorlinks=true, linkcolor= RawSienna, citecolor = RawSienna, filecolor = RawSienna, urlcolor = RawSienna, hypertexnames = true, backref = page]{hyperref}
\usepackage[]{natbib} 
\bibpunct[:]{(}{)}{,}{a}{}{,}
\geometry{left = 1.0in,right = 1.0in,top = 1.0in,bottom = 1.0in}
\usepackage[english]{babel}
\usepackage{float}
\usepackage{caption}
\usepackage{subcaption}
\usepackage{tikz}
\usepackage{booktabs}
\usepackage{pdfpages}
\usepackage{threeparttable}
\usepackage{framed}
\usepackage{comment}
\usepackage{lscape}
\usepackage{bm}
\setstretch{1.4}
%\usepackage[tablesfirst,nolists]{endfloat}

\usepackage[T1]{fontenc}
\usepackage{mlmodern}  % 太いComputer Modern
% MLmodernのバグを修正: cf. https://tex.stackexchange.com/questions/646333/size-of-integral-symbol-in-section-header-with-mlmodern
\DeclareFontFamily{OMX}{mlmex}{}
\DeclareFontShape{OMX}{mlmex}{m}{n}{<->mlmex10}{} 
\usepackage{tgtermes} % 数式以外の欧文をTXフォントで上書き

% Theorem numbering setup
\newtheorem{theorem}{Theorem}
\newtheorem{assumption}{Assumption}
\newtheorem{lemma}{Lemma}
\newtheorem{definition}{Definition}
\newtheorem{proposition}{Proposition}
\newtheorem{claim}{Claim}
\newtheorem{corollary}{Corollary}
\newtheorem{example}{Example}
\newtheorem{result}{Result}
\newtheorem{remark}{Remark}



%theorem環境(番号なし)
\newtheorem*{theorem*}{Theorem}
\newtheorem*{assumption*}{Assumption}
\newtheorem*{lemma*}{Lemma}
\newtheorem*{definition*}{Definition}
\newtheorem*{proposition*}{Proposition}
\newtheorem*{claim*}{Claim}
\newtheorem*{corollary*}{Corollary}
\newtheorem*{example*}{Example}
\newtheorem*{result*}{Result}
\newtheorem*{remark*}{Remark}





\DeclareMathOperator{\rank}{rank}
\hyphenation{non-iden-ti-fi-ca-tion}

% Appendix theorem numbering
\makeatletter
\newcommand{\appendixsection}{%
  \setcounter{section}{0}%
  \renewcommand{\thesection}{\Alph{section}}%
  \renewcommand{\thetheorem}{\Alph{section}.\arabic{theorem}}%
  \renewcommand{\thelemma}{\Alph{section}.\arabic{lemma}}%
  \renewcommand{\theproposition}{\Alph{section}.\arabic{proposition}}%
  \renewcommand{\thecorollary}{\Alph{section}.\arabic{corollary}}%
  \renewcommand{\thedefinition}{\Alph{section}.\arabic{definition}}%
}
\makeatother




\title{Revisiting the Identification of the Conduct Parameter in Homogeneous Goods Markets}
\author{Yuri Matsumura\thanks{Department of Economics, Rice University. Email: \href{mailto:yuri.matsumura23@gmail.com}{yuri.matsumura23@gmail.com}} \and Suguru Otani \thanks{Market Design Center, Department of Economics, University of Tokyo. Email: \href{mailto:suguru.otani@e.u-tokyo.ac.jp}{suguru.otani@e.u-tokyo.ac.jp}
}}

\begin{document}

\section{Model}


Consider homogeneous product markets.
The inverse demand function is $P(Q, X^{d})$ and individual marginal cost function is $mc_i(q_i, X^{s}_{i})$ where $q_i$ is the quantity produced by firm $i$, $Q=\sum_{i=1}^{N}q_i$ is the aggregate quantity, $X^{d}$ is the market level demand shifter, and $X^{s}_i$ is the cost shifter which can potentially include firm level cost shifter such as the distance from a factory to a sales store.

Then the equilibrium condition is given as
\begin{align}
    P(Q, X^{d}) + \theta q_i \frac{\partial P}{\partial q_i}(Q, X^{d}) = mc_i(q_i, X^{s}_i)
\end{align}
Note that as $Q$ is the aggregate quantity, we have
\begin{align}
    \frac{\partial P}{\partial q_i}(Q, X^{d}) = \frac{\partial P}{\partial Q}(Q, X^{d})\frac{\partial Q}{\partial q_i} = \frac{\partial P}{\partial Q}(Q, X^{d}).
\end{align}
Thus, the equilibrium condition is
\begin{align}
    P(Q, X^{d}) + \theta q_i \frac{\partial P}{\partial Q}(Q, X^{d}) = mc_i(q_i, X^{s}_i)
\end{align}

The definition of the non-identification is the following:
\begin{definition}
    
\end{definition}
 

This implies that observational equivalence holds for any value of demand shifters and cost shifters: $h_i^{*}(X^{d}, X^{s}_i) = h_i(X^{d}, X^{s}_i)$ for any $X^{d}$ and $X^{s}_i$.

The implicit fu









\end{document}